\documentclass[UTF8]{ctexart}
\usepackage{xeCJK}
\usepackage{geometry}
\geometry{left=2cm,right=2cm,top=2.5cm,bottom=2.5cm}
\setlength{\headheight}{12.7pt}

\usepackage{titlesec}
\titleformat{\section}
  {\normalfont\large\bfseries}
  {\thesection}
  {1em}
  {}
\titlespacing*{\section}
  {0pt}
  {3.5ex plus 1ex minus .2ex}
  {2.3ex plus .2ex}

\usepackage{amsmath}
\usepackage{amssymb}
\usepackage{amsthm}
\usepackage{enumitem}
\setlist[enumerate]{itemsep=0pt,topsep=0pt,parsep=0pt,leftmargin=0pt,itemindent=2.5em}
\setlist[itemize]{itemsep=0pt,topsep=0pt,parsep=0pt,leftmargin=0pt,itemindent=2.5em}
\usepackage{fancyhdr}
\pagestyle{fancy}
\fancyhf{}
\usepackage{graphicx}
\usepackage{hyperref}
\usepackage{indentfirst}
\usepackage{lmodern}


\begin{document}
\setlength{\abovedisplayskip}{1pt}
\setlength{\belowdisplayskip}{1pt}

% 久违的轻松.
% 2026.04.02

\section{二元关系的定义}

对任意的\href{01 集合的简单运算#nameddest=sec:定义2.2}{有序对}
$p=(x,y)$, 这里记$x=p_L,y=p_R$, 记$p^\mathrm{C}=(y,x)$.

\vspace{10pt}

\hypertarget{sec:定义1.1}{}
\noindent\textbf{定义1.1 (二元关系)}

给定集合$r$, 如果任意的$p\in r$都是有序对, 那么称$r$是\textbf{二元关系}.

进一步, 如果$r\subset a\times b$, 称$r$是从$a$到$b$的二元关系;
如果$r\subset a\times a$, 称$r$是$a$上的二元关系.

\vspace{10pt}

\hypertarget{sec:定义1.2}{}
\noindent\textbf{定义1.2 (定义域与值域)}

给定二元关系$r$, 定义
\[
    \mathrm{Dom}\,r\overset{\mathrm{def}}{=}\{p_L\mid p\in r\},\quad
    \mathrm{Ran}\,r\overset{\mathrm{def}}{=}\{p_R\mid p\in r\},
\]
分别称为$r$的\textbf{定义域}和\textbf{值域}.

显然对任意的$x$, $x\in\mathrm{Dom}\,r$当且仅当$\exists y:(x,y)\in r$;

对任意的$y$, $y\in\mathrm{Ran}\,r$当且仅当$\exists x:(x,y)\in r$.

\vspace{10pt}

\hypertarget{sec:定理1.1}{}
\noindent\textbf{定理1.1}

给定二元关系$r$,
\begin{enumerate}
    \item $r\subset\mathrm{Dom}\,r\times\mathrm{Ran}\,r$,
    \item 如果$r\subset a\times b$,
    那么$\mathrm{Dom}\,r\subset a$, $\mathrm{Ran}\,r\subset b$.
\end{enumerate}

\noindent{\kaishu 证明.}

\begin{enumerate}
    \item 对任意的$p\in r$, $p_L\in\mathrm{Dom}\,r$且$p_R\in\mathrm{Ran}\,r$,
    所以$p=(p_L,p_R)\in\mathrm{Dom}\,r\times\mathrm{Ran}\,r$.
    \item 假设$r\subset a\times b$. 对任意的$x\in\mathrm{Dom}\,r$,
    存在$y$使得$(x,y)\in r$, 从而$x\in a$;

    \hspace{2.17em}
    对任意的$y\in\mathrm{Ran}\,r$, 存在$x$使得$(x,y)\in r$, 从而$y\in b$.
    \hfill $\qedsymbol$
\end{enumerate}

\vspace{10pt}

\hypertarget{sec:定理1.2}{}
\noindent\textbf{定理1.2}

给定二元关系$r,s$,
\begin{enumerate}
    \item 如果$r\subset s$, 那么$\mathrm{Dom}\,r\subset\mathrm{Dom}\,s$,
    $\mathrm{Ran}\,r\subset\mathrm{Ran}\,s$,
    \item $\mathrm{Dom}(r\cup s)=\mathrm{Dom}\,r\cup\mathrm{Dom}\,s$,
    $\mathrm{Ran}(r\cup s)=\mathrm{Ran}\,r\cup\mathrm{Ran}\,s$.
    \item $\mathrm{Dom}(r\cap s)\subset\mathrm{Dom}\,r\cap\mathrm{Dom}\,s$,
    $\mathrm{Ran}(r\cap s)\subset\mathrm{Ran}\,r\cap\mathrm{Ran}\,s$.
    \item $\mathrm{Dom}(r\setminus s)\supset
    \mathrm{Dom}\,r\setminus\mathrm{Dom}\,s$,
    $\mathrm{Ran}(r\setminus s)\supset\mathrm{Ran}\,r\setminus\mathrm{Ran}\,s$.
\end{enumerate}

\vspace{10pt}

\hypertarget{sec:定义1.3}{}
\noindent\textbf{定义1.3 (逆关系)}

给定二元关系$r$, 定义它的\textbf{逆关系}为
$r^{-1}\overset{\mathrm{def}}{=}\{p^\mathrm{C}\mid p\in r\}$.

\vspace{10pt}

\hypertarget{sec:定理1.3}{}
\noindent\textbf{定理1.3}

给定二元关系$r,s$,
\begin{enumerate}
    \item $(r^{-1})^{-1}=r$,
    \item $\mathrm{Dom}\,r^{-1}=\mathrm{Ran}\,r$,
    $\mathrm{Ran}\,r^{-1}=\mathrm{Dom}\,r$,
    \item 如果$r\subset s$, 那么$r^{-1}\subset s^{-1}$,
    \item $(r\cup s)^{-1}=r^{-1}\cup s^{-1}$, $(r\cap s)^{-1}=r^{-1}\cap s^{-1}$,
    $(r\setminus s)^{-1}=r^{-1}\setminus s^{-1}$.
\end{enumerate}





\newpage

\section{二元关系的限制}

\hypertarget{sec:定义2.1}{}
\noindent\textbf{定义2.1 (二元关系的限制)}

给定二元关系$r$和集合$a$, 定义$r$在$a$上的\textbf{限制}
\[
    r\!\restriction_a\overset{\mathrm{def}}{=}\{p\in r\mid p_L\in a\}.
\]
显然$r\!\restriction_a$是$r$的子集, 也是一个二元关系.

\vspace{10pt}

\hypertarget{sec:定理2.1}{}
\noindent\textbf{定理2.1}

给定二元关系$r$和集合$a,b$,
\begin{enumerate}
    \item 如果$a\subset b$, 那么$r\!\restriction_a\,\subset r\!\restriction_b$,
    \item $r\!\restriction_{a\cup b}\,=r\!\restriction_a\cup\,r\!\restriction_b$,
    $r\!\restriction_{a\cap b}\,=r\!\restriction_a\cap\,r\!\restriction_b$,
    $r\!\restriction_{a\setminus b}\,=r\!\restriction_a\setminus\,r\!\restriction_b$.
\end{enumerate}

\vspace{10pt}

\hypertarget{sec:定义2.2}{}
\noindent\textbf{定义2.2 (像与原像)}

给定二元关系$r$和集合$a$, 定义
\[
    r[a]\overset{\mathrm{def}}{=}\mathrm{Ran}\,r\!\restriction_a,
\]
称为$a$在$r$下的\textbf{像}.

进一步, 再给定集合$b$, 称$r^{-1}[b]$为$b$在$r$下的\textbf{原像}.

\vspace{10pt}

\hypertarget{sec:定理2.2}{}
\noindent\textbf{定理2.2}

给定二元关系$r$和集合$a$,
\begin{enumerate}
    \item $\mathrm{Dom}\,r\!\restriction_a=a\cap\mathrm{Dom}\,r$,
    \item 对任意的$y$, $y\in r[a]$当且仅当$\exists x\in a:(x,y)\in r$.
\end{enumerate}

\noindent{\kaishu 证明.}

\begin{enumerate}
    \item 对任意的$x$, $x\in\mathrm{Dom}\,r\!\restriction_a\iff
    \exists y:(x,y)\in r\!\restriction_a\iff
    \exists y:((x,y)\in r\land x\in a)\iff x\in\mathrm{Dom}\,r\land x\in a$.
    \item 对任意的$y$, $y\in r[a]\iff y\in\mathrm{Ran}\,r\!\restriction_a\iff
    \exists x:(x,y)\in r\!\restriction_a\iff\exists x:((x,y)\in r\land x\in a)$.
    \hfill $\qedsymbol$
\end{enumerate}

\vspace{10pt}

\hypertarget{sec:定理2.3}{}
\noindent\textbf{定理2.3}

给定二元关系$r$和集合$a,b$,
\begin{enumerate}
    \item 如果$a\subset b$, 那么$r[a]\subset r[b]$,
    \item $r[a\cup b]=r[a]\cup r[b]$,
    \item $r[a\cap b]\subset r[a]\cap r[b]$,
    \item $r[a\setminus b]\supset r[a]\setminus r[b]$,
    \item $r^{-1}[r[a]]\supset a\cap\mathrm{Dom}\,r$.
\end{enumerate}





\newpage

\section{二元关系的复合}

\hypertarget{sec:定义3.1}{}
\noindent\textbf{定义3.1 (二元关系的复合)}

给定二元关系$r$和$s$, 定义它们的\textbf{复合}
\[
    s\circ r\overset{\mathrm{def}}{=}
    \{(x,z)\in\mathrm{Dom}\,r\times\mathrm{Ran}\,s\mid
    \exists y:((x,y)\in r\land(y,z)\in s)\}.
\]
显然$s\circ r$是从$\mathrm{Dom}\,r$到$\mathrm{Ran}\,s$的二元关系.

显然对任意的$x,z$, $(x,z)\in s\circ r$当且仅当存在$y$使得
$(x,y)\in r$且$(y,z)\in s$.

\vspace{10pt}

\hypertarget{sec:定理3.1}{}
\noindent\textbf{定理3.1}

给定二元关系$r,r',s,s'$和集合$a$,
\begin{enumerate}
    \item 如果$r\subset r',s\subset s'$, 那么$s\circ r\subset s'\circ r'$,
    \item $(s\circ r)^{-1}=r^{-1}\circ s^{-1}$,
    \item $\mathrm{Ran}(s\circ r)=s[\,\mathrm{Ran}\,r\,]$,
    $\mathrm{Dom}(s\circ r)=r^{-1}[\,\mathrm{Dom}\,s\,]$,
    \item $(s\circ r)\!\restriction_a=s\circ r\!\restriction_a$,
    $(s\circ r)[a]=s[r[a]]$.
\end{enumerate}

\noindent{\kaishu 证明.}

\begin{enumerate}
    \item 如果$r\subset r',s\subset s'$, 那么对任意的$(x,z)\in s\circ r$,
    存在$y$使得
    \[
        (x,y)\in r\land(y,z)\in s\quad\implies\quad
        (x,y)\in r'\land(y,z)\in s',
    \]
    所以$(x,z)\in s'\circ r'$.

    \item 对任意的$(z,x)\in(s\circ r)^{-1}$, 有$(x,z)\in s\circ r$,
    即存在$y$使得
    \[
        (x,y)\in r\land(y,z)\in s\quad\implies\quad
        (z,y)\in s^{-1}\land(y,x)\in r^{-1},
    \]
    所以$(z,x)\in r^{-1}\circ s^{-1}$,
    即$(s\circ r)^{-1}\subset r^{-1}\circ s^{-1}$.

    \hspace{2.17em}
    据此有$(r^{-1}\circ s^{-1})^{-1}\subset s\circ r$,
    于是$r^{-1}\circ s^{-1}\subset(s\circ r)^{-1}$.

    \item 对任意的$z$,
    \[\begin{aligned}
        z\in\mathrm{Ran}(s\circ r)\iff&\exists x:(x,z)\in s\circ r\iff
        \exists x\,\exists y:((x,y)\in r\land(y,z)\in s)\\[-3pt]
        \iff&\exists y\in\mathrm{Ran}\,r:(y,z)\in s\iff
        z\in s[\,\mathrm{Ran}\,r\,].\\[1pt]
    \end{aligned}\]
    进一步, $\mathrm{Dom}(s\circ r)=\mathrm{Ran}(s\circ r)^{-1}=
    \mathrm{Ran}(r^{-1}\circ s^{-1})=r^{-1}[\,\mathrm{Ran}\,s^{-1}\,]=
    r^{-1}[\,\mathrm{Dom}\,s\,]$.

    \item 对任意的$(x,z)$, 有
    \[
        (x,z)\in(s\circ r)\!\restriction_a\iff
        x\in a\land\exists y:(x,y)\in r\land(y,z)\in s\iff
        (x,z)\in s\circ r\!\restriction_a.
    \]
    于是$(s\circ r)\!\restriction_a=s\circ r\!\restriction_a$,
    从而$(s\circ r)[a]=s[r[a]]$. \hfill $\qedsymbol$
\end{enumerate}

\vspace{10pt}

\hypertarget{sec:定理3.2}{}
\noindent\textbf{定理3.2 (复合的结合律)}

给定二元关系$r,s,t$, $(t\circ s)\circ r=t\circ(s\circ r)$.

\noindent{\kaishu 证明.}

对任意的$x,w$,
\begin{equation*}\begin{aligned}
    &(x,w)\in(t\circ s)\circ r\\
    \iff&\exists y:\quad(x,y)\in r\land(y,w)\in t\circ s\\
    \iff&\exists y\,\exists z:\quad(x,y)\in r\land(y,z)\in s\land(z,w)\in t\\
    \iff&\exists z:\quad(x,z)\in s\circ r\land(z,w)\in t\\
    \iff&(x,w)\in t\circ(s\circ r).
\end{aligned}\tag*{\qedsymbol}\end{equation*}

\end{document}